(1) 固层满足 \(T_t=\alpha T_{xx}\)、\(T(0,t)=T_0\)、\(T(s(t),t)=T_m\)。令 \(\eta=x/(2\sqrt{\alpha t})\)、\(s=2\lambda\sqrt{\alpha t}\)，自相似常微分方程是 \(f''+2\eta f'=0\)。两次积分并施加边界值得
\[
\boxed{T(x,t)=T_0+(T_m-T_0)\frac{\operatorname{erf}(\eta)}{\operatorname{erf}(\lambda)}}.
\]
对每一固定 \(t>0\)，此场严格满足导热微分方程；\(\lambda\) 必须由相界能量条件决定，不能随意作为额外自由参数。

(2) Stefan条件是 \(\rho L\dot s=kT_x(s^-,t)\)。使用 \(\dot s=\lambda\sqrt{\alpha/t}\) 和
\(T_x(s^-)=\Delta T\,e^{-\lambda^2}/(\sqrt{\pi\alpha t}\operatorname{erf}\lambda)\)，得到
\[
\boxed{\lambda e^{\lambda^2}\operatorname{erf}\lambda=\frac{c_p\Delta T}{L\sqrt\pi}=\frac{\mathrm{Ste}}{\sqrt\pi}}.
\]
函数 \(F(\lambda)=\lambda e^{\lambda^2}\operatorname{erf}\lambda\) 在 \(\lambda>0\) 上导数严格正，且由0增长至无穷，所以每个正Stefan数有且仅有一个正根。故 \(s\propto\sqrt t\)，而不是界面匀速前进。

(3) 壁面向外热流大小为 \(q_0(t)=kT_x(0,t)=k\Delta T/(\sqrt{\pi\alpha t}\operatorname{erf}\lambda)\)，累计移热
\(Q_0(t)=2k\Delta T\sqrt t/(\sqrt{\pi\alpha}\operatorname{erf}\lambda)\)（按单位面积）。固层相对于原来 \(T_m\) 液体的能量亏损是
\[
Q_{\rm stored}=\rho Ls+\rho c_p\int_0^s(T_m-T)\,dx.
\]
误差函数积分 \(\int_0^\lambda\operatorname{erf}\eta\,d\eta=\lambda\operatorname{erf}\lambda+(e^{-\lambda^2}-1)/\sqrt\pi\)，故显热项为
\(2\rho c_p\Delta T\sqrt{\alpha t}(1-e^{-\lambda^2})/(\sqrt\pi\operatorname{erf}\lambda)\)。再用上问Stefan条件可逐项化简至 \(Q_{\rm stored}=Q_0\)，这同时检查了壁面热流和界面潜热的正号。

(4) 小 \(\lambda\) 时 \(\operatorname{erf}\lambda=(2/\sqrt\pi)(\lambda-\lambda^3/3+\cdots)\)，所以 \(2\lambda^2+O(\lambda^4)=\mathrm{Ste}\)，得到 \(\lambda\sim\sqrt{\mathrm{Ste}/2}\) 与 \(s\sim\sqrt{2\alpha\mathrm{Ste}\,t}\)。大Stefan数 \(\lambda\gg1\)，\(\operatorname{erf}\lambda\to1\)，故 \(\lambda^2+\log\lambda\sim\log(\mathrm{Ste}/\sqrt\pi)\)，首阶 \(\lambda\sim\sqrt{\log\mathrm{Ste}}\)。\(q_0\sim t^{-1/2}\) 虽在初瞬发散，但其时间积分有限；恒定温度边界、无液侧显热和忽略成核等理想化不能外推到微观 \(t=0\)。
